% =============================================================================
% Ready-to-paste LaTeX: the waveform-domain representation (wav2vec 2.0) against the
% spectrogram-based ones. Answers the supervisor request "try a waveform extractor,
% not just spectrogram".
%
% Numbers: emotion regression (tab:waveform_stage1) from
% experiments/features/exp_waveform_vs_spectrogram.py (results/waveform_vs_spectrogram.json);
% the layer sweep from experiments/features/exp_w2v_layer_sweep.py
% (results/w2v_layer_sweep.json); genre classification (tab:waveform_stage2) from the
% corrected run, experiments/evaluation/exp_cv_corrected.py (results/cv_corrected.json),
% where all six representations share the headline 10x5 folds. Seed 42 throughout.
% Do NOT reintroduce the earlier 5x5 per-fold genre figures (.391 / .376 / .342 / .300);
% they are superseded (progress log section 21).
%
% Placement: the feature-extraction part of the Methods chapter for the first
% subsection, and the Results chapter for the two tables. The layer-selection
% subsection can be moved to an appendix if space is tight, but it should be kept
% somewhere -- it is the control that makes the negative result defensible.
%
% CITATION NOTE: uses baevski2020wav2vec and pons2019musicnn. Both are now in
% bib/library.bib on Overleaf; the wav2vec entry is reproduced here for reference:
%
%   @inproceedings{baevski2020wav2vec,
%     title     = {wav2vec 2.0: A Framework for Self-Supervised Learning of Speech
%                  Representations},
%     author    = {Baevski, Alexei and Zhou, Yuhao and Mohamed, Abdelrahman and
%                  Auli, Michael},
%     booktitle = {Advances in Neural Information Processing Systems (NeurIPS)},
%     volume    = {33}, pages = {12449--12460}, year = {2020}
%   }
% =============================================================================

\section{Waveform versus Spectrogram Representations}
\label{sec:waveform}

Each of the audio representations examined so far begins by converting the signal into a
time--frequency image. AST and CLAP operate on log-mel spectrograms, VGGish and MusiCNN on
log-mel patches, and the hand-crafted baseline on descriptors derived from the short-time
Fourier transform. The comparison between emotion features and learned audio embeddings has
therefore only been conducted against \emph{spectrogram} front-ends, which leaves open the
objection that the limitation lies in the mel representation rather than in the idea of
classifying genre directly from audio.

To close that gap, wav2vec 2.0 \citep{baevski2020wav2vec} is added as a sixth
representation. Its feature encoder is a stack of one-dimensional convolutions applied
directly to the raw 16\,kHz sample sequence; no spectral transform occurs at any stage.
The self-supervised base checkpoint is used rather than an ASR fine-tune, whose
representations are specialised toward phonetic content. Clip length
(10.24\,s), temporal mean pooling and output dimensionality (768) are identical to AST, so
that the two differ only in their front-end.

One confound must be stated plainly, because it constrains what the comparison can show.
wav2vec 2.0 was pretrained on read speech, whereas AST and VGGish were pretrained on
AudioSet, which contains music. The comparison therefore varies the input domain and the
pretraining corpus simultaneously, and a poor result for wav2vec 2.0 cannot be attributed
to the waveform front-end alone.

MusiCNN \citep{pons2019musicnn} is included for the complementary reason. It is pretrained
on Million Song Dataset tagging, so of all the representations considered it matches the
material being analysed most closely, and it is the natural candidate if domain-specific
pretraining is what the task needs. Together the six span four pretraining regimes ---
general audio events (AST, VGGish), audio--text (CLAP), music tagging (MusiCNN) and speech
(wav2vec 2.0), against a hand-crafted descriptor baseline --- so that the comparison
against emotion features is not a comparison against a single, possibly badly chosen,
embedding.

\subsection{Selecting the Layer to Pool}

A self-supervised model's final layers specialise toward its own pretraining objective, so
pooling the last hidden state --- the obvious default --- is not necessarily the best
choice for a task the model was never trained on. All thirteen hidden states were
therefore evaluated on the emotion-regression task under the protocol used elsewhere
(random-forest regression, film-grouped cross-validation, mean $R^2$ across the eight
emotions). Performance peaks at the output of the second transformer block
($R^2 = 0.324$) and falls away thereafter --- not strictly monotonically, since the tenth
block recovers to $0.187$ --- reaching $R^2 = 0.161$ at the final (twelfth) block, with
the convolutional encoder output alone already reaching $0.292$.

The choice of layer is thus worth a factor of two --- larger than the effect this thesis
sets out to measure --- and reporting the final layer by default would have understated
the waveform representation substantially. The second block is used in all results below.
It is
selected on the emotion-regression task and then held fixed; it is not tuned against any
genre result.

\subsection{Results}

Table~\ref{tab:waveform_stage1} reports emotion regression, the metric by which feature
extractors are ranked (Section~\ref{sec:evaluation_protocol}). The waveform representation
recovers markedly less of the emotion-rating variance than any spectrogram-based
alternative, and less even than the 103-dimensional hand-crafted baseline. The deficit is
uniform across the eight emotions rather than concentrated in any one of them.

The music-pretrained representation does not lead either. MusiCNN reaches $R^2 = 0.541$,
clearly above the hand-crafted descriptors and far above the waveform model, but below
all three general-purpose spectrogram embeddings. Domain-matched pretraining is thus
worth less here than might be expected --- the ranking is driven by how much
emotion-relevant variance a representation exposes, not by how closely its pretraining
corpus resembles film music.

\begin{table}[t]
  \centering
  \caption{Emotion regression by audio representation (random forest, film-grouped
    cross-validation, mean $R^2$ over the eight emotions, $n=360$ clips). The ceiling
    column expresses each score as a fraction of the $R^2 = 0.897$ implied by the
    inter-panel reliability of the ratings.}
  \label{tab:waveform_stage1}
  \begin{tabular}{llrrrr}
    \hline
    Representation & Input domain & Dim. & $R^2$ & RMSE & \% ceiling \\
    \hline
    CLAP           & log-mel spectrogram & 512 & .561 & 1.22 & 62 \\
    AST            & log-mel spectrogram & 768 & .560 & 1.22 & 62 \\
    VGGish         & log-mel spectrogram & 128 & .558 & 1.23 & 62 \\
    MusiCNN        & log-mel spectrogram & 200 & .541 & 1.25 & 60 \\
    Hand-crafted MIR & STFT descriptors  & 103 & .490 & 1.31 & 55 \\
    \textbf{wav2vec 2.0} & \textbf{raw waveform} & 768 & \textbf{.323} & 1.52 & 36 \\
    \hline
  \end{tabular}
\end{table}

Table~\ref{tab:waveform_stage2} reports genre classification on the five-genre subset. Two
results stand out, and the second is the more important of the two.

\begin{table}[t]
  \centering
  \caption{Genre classification by representation (five-genre subset, $10\times5$
    repeated film-grouped cross-validation, nested-CV-tuned regularisation; macro-$F_1$ on
    the pooled out-of-fold predictions with a 95\,\% interval from a bootstrap over films).}
  \label{tab:waveform_stage2}
  \begin{tabular}{lrl}
    \hline
    Approach & Macro-$F_1$ & 95\,\% CI \\
    \hline
    Emotion (11), ground-truth ratings                       & .428 & [.371, .476] \\
    \textbf{wav2vec 2.0 $\rightarrow$ predicted emotion (11)} & \textbf{.419} & [.367, .462] \\
    VGGish-128, direct                                        & .377 & [.325, .414] \\
    AST-768, direct                                           & .371 & [.326, .398] \\
    Hand-crafted MIR, direct                                  & .363 & [.318, .398] \\
    MusiCNN-200, direct                                       & .363 & [.311, .401] \\
    CLAP-512, direct                                          & .361 & [.313, .400] \\
    \textbf{wav2vec 2.0, direct}                              & \textbf{.326} & [.288, .356] \\
    \hline
  \end{tabular}
\end{table}

First, used directly, the waveform representation is the weakest of the six
($0.326$). Its shortfall against the other five is small and only partly significant,
however: $+0.034$ to $+0.049$, significant under the film bootstrap against VGGish, AST
and the hand-crafted descriptors ($p = 0.014$--$0.040$) but not against CLAP or MusiCNN
($p \approx 0.10$), and against none of them under the stricter Nadeau--Bengio test
($p = 0.17$--$0.38$). This is itself informative: the genre task discriminates far less
sharply between representations than emotion regression does, where the same six differ
by a factor of $1.7$. This supports
both the choice of $R^2$ as the metric for ranking feature extractors and the conclusion
of Section~\ref{sec:error_analysis} that genre classification on this corpus operates near
a noise ceiling that a better representation cannot lift.

Second, and centrally, \emph{the emotion bottleneck rescues the weakest representation}.
Routed through predicted emotions, wav2vec 2.0 rises from last place to second, reaching
$0.419$ --- above every directly used embedding, including the spectrogram-based ones, and
statistically indistinguishable from the ground-truth emotion ceiling
($-0.009$, $p = 0.46$). The improvement over its own direct use is $+0.091$, significant
under both tests (film bootstrap $p < 0.001$, Nadeau--Bengio $p = 0.012$), and it holds
in \emph{all} ten repeats of the cross-validation. The benefit of the
emotion intermediate is therefore not a property of one particular embedding: it is
largest precisely where the underlying representation is weakest, which is the behaviour
expected of a bottleneck that discards representation-specific detail and retains only
perceptually meaningful structure.
